\section{極限と連続}\label{節:関数極限連続}
数列では自然数 $n$ を大きくした.関数では,実数 $x$ をある点 $a$ に近づける.\,$x\ne a$ のまま $x$ を $a$ に十分近づければ,\,$f(x)$ を $L$ にどれだけでも近づけられるとき,
\[
 \lim_{x\to a}f(x)=L
\]
と書く.正確には,任意の $\varepsilon>0$ に対して,ある $\delta>0$ を選び,\,$0<|x-a|<\delta$ なら $|f(x)-L|<\varepsilon$ となることである.関数の定義域に属する $x$ だけを考える.

$x>a$ の側から近づける右極限を $\lim_{x\to a+}f(x)$,\,$x<a$ の側から近づける左極限を $\lim_{x\to a-}f(x)$ と書く.両側で近づけられる点では,左右の極限が同じ値に一致して初めて両側の極限が存在する.たとえば $\frac{|x|}{x}$ は $x\to0+$ で $1$,\,$x\to0-$ で $-1$ だから,\,$x\to0$ の極限は存在しない.

\begin{definition}{連続性}
$f$ が $x=a$ で定義され,\,$\lim_{x\to a}f(x)=f(a)$ となるとき,\,$f$ は $a$ で\textbf{連続}であるという.閉区間の端点では,区間の内側からの片側極限を使う.
\end{definition}

極限は $x=a$ 以外の近くの値から決まるのに対し,連続性はその極限と実際の値 $f(a)$ が一致するかどうかも問う.多項式,分母が $0$ でない有理関数,および三角関数・指数関数・対数関数は,それぞれの定義域で連続である.和・積・商・合成による連続性の保存も基本性質として用いる.

\begin{example}{値を補って連続にする}
$x\ne2$ に対して $f(x)=\frac{x^2-4}{x-2}$ と定める.\,$f(2)$ をどのように定めれば,\,$x=2$ でも連続になるか.
\end{example}
$x\ne2$ の範囲では約分でき,
\[
 f(x)=\frac{(x-2)(x+2)}{x-2}=x+2.
\]
したがって $\lim_{x\to2}f(x)=4$ である.\,$f(2)=4$ と定めれば極限と関数値が一致し,連続になる.約分後の $x+2$ は $x=2$ でも値をもつが,元の分数はそのままでは定義されない.式が等しい範囲と,関数の定義域を分けて考える必要がある.

後で微分の定理を証明するために,連続関数について次の二つの存在定理を使う.

\begin{theorem}{連続関数の存在定理}\label{定理:連続関数存在}
$f$ が閉区間 $[a,b]$ 上で連続なら,次が成り立つ.
\begin{enumerate}[label=(\arabic*),leftmargin=*]
 \item $f$ は区間内で最大値と最小値をとる.
 \item $f(a)$ と $f(b)$ の間にある任意の値 $k$ に対して,\,$f(c)=k$ となる $c\in[a,b]$ が存在する.これを\textbf{中間値の定理}という.
\end{enumerate}
\end{theorem}

これらの証明も実数の連続性に基づくため,本書では採用する事実とする.条件を外すと結論は崩れる.たとえば $f(x)=x$ は開区間 $(0,1)$ では最大値も最小値もとらない.端の値 $0,1$ に近づくことはできるが,区間の中には到達する点がない.閉区間という条件は,単なる表記上の飾りではない.


